\begin{table}

\caption{Matched text with before and after context for selected reuse examples. Portions with highest similarity highlighted. DOIs of original publications given.}\label{tab:example-cases}
\centering
\begin{tabularx}{\textwidth}{@{}XX@{}}
\toprule
\multicolumn{1}{c}{\textsf{\textbf{Publication A}}} & \multicolumn{1}{c}{\textsf{\textbf{Publication B}}} \\ \midrule

[\dots] children than men in lower classes nor experienced the same downward mobility. Clark writes, \ctext[RGB]{255,214,75}{``Thus we may speculate that England's advantage lay in the rapid cultural, and potentially also genetic, diffusion of the values of the economically successful throughout society in the years 1200-1800'' (p. 271).} It does not serve the book well to dwell on the rebuttal of its evolutionary [\dots] \newline {\scriptsize\color{gray}{10.1111/j.1475-4991.2009.00354.x}}
& [\dots] Impact on the standard of living, but eventually they led to the end of the long Malthusian era. \ctext[RGB]{255,214,75}{``Thus we may speculate that England's advantage lay in the rapid cultural, and potentially also genetic, diffusion of the values of the economically successful throughout society in the years 1200-1800'' (p. 271).} Finally, Clark considers the great divergence among today's economies. Because [\dots] \newline {\scriptsize\color{gray}{10.1111/j.1468-0289.2008.00432\_10.x}}
\\

\midrule

[\dots] static approximation of the ColourSinglet Model, we have seen that the production amplitude \ctext[RGB]{255,214,75}{receives contributions from two different cuts. The first one in its static limit gives the colour-singlet mechanism. The second one has not been considered so far. We treat it in a gauge-invariant manner by introducing necessary new 4-point vertices, suggestive of the colour-octet mechanism.} This new contribution can be as large as the colour-singlet mechanism at high \newline {\scriptsize\color{gray}{10.1063/1.2122163}}
& [\dots] and we show that the lowest-order mechanism for heavy-quarkonium production \ctext[RGB]{255,214,75}{receives in general contributions from two different cuts. The first one corresponds to the usual colour-singlet mechanism. The second one has not been considered so far. We treat it in a gauge-invariant manner, and introduce new 4-point vertices, suggestive of the colour-octet mechanism}. These new objects enable us to go beyond the static approximation. We show that the contribution of [\dots] \newline {\scriptsize\color{gray}{10.1016/j.physletb.2005.11.073}}
\\

\midrule

[\dots] \ctext[RGB]{255,214,75}{The authors declare that they have no competing interests.} MS and MMD drafted and wrote the manuscript. MS, MMD and KCRP participated in the care of the patient and interpretation of the investigations. \ctext[RGB]{255,214,75}{All authors read and approved the final manuscript.} [\dots] \newline {\scriptsize\color{gray}{10.1186/1757-1626-2-147}}
& [\dots] \ctext[RGB]{255,214,75}{The authors declare that they have no competing interests.} Authors' contributions: AJB is the chief investigator of the CASP trial, he is responsible for the conduct of the study and the clinical supervision and training of the therapists and he wrote the first draft of this paper. CS drafted an earlier version of the grant proposal and collaborated with AJB, MT, RR \& PH to produce the successful proposal. CS also provided clinical supervision and therapist training. LS edited and revised the proposal and is responsible for the day to day running of the trial. \ctext[RGB]{255,214,75}{All authors read and approved the final manuscript.} [\dots] \newline {\scriptsize\color{gray}{10.1186/1471-244X-13-199}} \\

\bottomrule
\end{tabularx}

\end{table}

